\documentclass[12pt]{article}
\usepackage[utf8]{inputenc}
\usepackage[T1]{fontenc}
\usepackage{hyperref}
\title{Yessan-Mayo Language Resources for K-12 Curriculum Development}
\author{Compiled by Claude Sonnet 5.5 (claude-sonnet-5-5) for Andr\'as Kornai}
\date{}
\begin{document}
\maketitle

\section{Introduction}
Yessan-Mayo (ISO 639-3 \texttt{yss}, Glottocode yess1239; also called Yesan, Mayo and, natively, Yamano) is a Sepik language of Papua New Guinea. The 2000 census gave about 2,000 speakers. The speakers live in Maio and Yessan villages, Ambunti Rural LLG, East Sepik Province. Wikitongues gives it EGIDS 5 (Developing). Joshua Project gives 4,100 for the wider Mayo-Yessan people. The search found no school materials, a dictionary or a teaching grammar.

\section{Primary Educational Resources}
No primers or school programmes were found.

\section{Linguistic and Reference Materials}
\begin{itemize}
\item Wikipedia, Yessan-Mayo language. It cites Foley (2018) for the phonology: \url{https://en.wikipedia.org/wiki/Yessan-Mayo_language}
\item WALS entry, which lists Foreman 1974, ``Grammar of Yessan-Mayo'', and Foreman and Marten 1973, ``Yessan-Mayo phonemes'', as sources: \url{https://wals.info/languoid/lect/wals_code_yes}
\item Max Planck EVA numeral-systems page, with data from Helen Marten (2010 and 2014): \url{https://lingweb.eva.mpg.de/channumerals/Yessan-Mayo.htm}
\end{itemize}

\section{Religious and Community Resources}
\begin{itemize}
\item The New Testament in Yamano, published by Lutheran Bible Translators and released under a Creative Commons licence (BY-NC-ND 4.0). This PDF copy of it is the one that appeared in the search: \url{https://nobanners.classicistranieri.com/biblepdf/yss-yamano/yss-yamano_all.pdf}
\item Joshua Project, Mayo-Yessan people: \url{https://joshuaproject.net/people_groups/19085}
\end{itemize}

\section{Digital and Audio Resources}
\begin{itemize}
\item Wikitongues page for yss: \url{https://wikitongues.org/languages/yss/}
\end{itemize}
Joshua Project reports audio recordings for the people group, but I saw no direct link.

\section{Curriculum Development Context}
The New Testament and the Foreman and Marten descriptions are the main written sources. The New Testament is a text of fixed content, and the licence forbids modifying it. A mother-tongue curriculum would need new primers and literacy materials built with the community, using the orthography of the existing translation.
\end{document}